% Pursuing Stacks (cotes 134-2 to 134-8): archive cote 134-7, site page 14 (PDF page 15), typescript page 456.
% typescript. Draft compared: ps6.tex:650-704.
% AI-transcribed by Claude Opus 5.5 (claude-opus-5-5, adaptive thinking, effort high), Message Batches API, batch ps3 (msgbatch_01XdFmuVVBAcgc7EFxdtCqsj), 2026-11-12; not yet checked by eye.
\scanpage{134-7/14}
\mspage{456}

\[ \Gamma_k : Ab_k \to Ab_k \qquad , \]
it is enough\deleted{\illegible{1}}\hand{,} \unsure{\hand{f}or} get\inserted{ting} it's value on a chain complex \xout{MxxxMxxxx} $M.'$ say ,
to replace \xout{Mx} $M.'$ by $M_*'$ and apply the functor $\Gamma_k$ componentwise,
without first having to take a projective resolution $M.$ of $M.'$ - some-%
thing rather absurd indeed ! Thus the statement made on page 454, that
when taking for $M.'$ the k-module \hand{$\pi$} placed in degree n and zero in all
other degrees, the corresponding $L_{pt}(X_*') = I_k(M_*')$ embodies the
value of the left derived functor $L\Gamma_k$ on $M.'$ , is visibly incorrect if
we do'nt assume moreover that \hand{$\pi$} is projective (flat presumably would be
enough ...). Otherwise, we should first replace \hand{$\pi$} by a projective
(or at any rate flat) resolution\hand{,} which we shift by n to get $M.$ , and \inserted{then}
take $\Gamma_k(M_*)$ to get the correct value of $L\Gamma_k(M.')$,

This convinces me that in the question as to whether the maps in (w)
are quasi-isomorphisms (the more crucial one of course being the first
of the two), \uline{we should assume moreover that the components of $X_*$ are}
\uline{described in terms of projective k-modules} $M_n$ , or at any rate k-modules
which are flat. Accordingly, we should make the same restriction on
the semisimplicial schematized model $X_*$ , in order for the description
we gave of ``k-linearization'' as $L_{pt}(X_*)$ (or $L(X_*)$, never mind which)
to be topologically meaningful. Very probably, in the whole schematiza-%
tion set-up, namely in the very definition we gave of $U(k)$ and $U(k)^\bullet$ ,
we should stick to the same restriction. If \xout{atxthexstart} I insisted
first (with some inner reluctance, I admit) on taking k-modules M
unrestricted, this was because I was thinking of M, more specifically
of the $M_i$'s occuring in the inductive ``dévissage'' of an object of
$U_m(k)$ (when thinking of the more refined version\deleted{s} of $U(k)$), as essen-%
tially the homotopy groups of the homotopy type we want to modelize,
or \inserted{rather, as} the components of the corresponding semisimplicial k-modules (denoted
$M_*'$ some minutes ago). I was still thinking of course,\deleted{\unsure{t}} be it implicit-%
ly, in terms of Postnikow dévissage, despite yesterday's remark that
to use such dévissage litterally may cause trouble (p.444). Thus the
feeling which gets into the fore now is that \uline{we should kind of forget}
\uline{Postnikov}, \uline{and} \deleted{\unsure{deal}} \inserted{work} \uline{with} \uline{\unsure{semisimplicial|Semisimplicial}} \uline{``schematic''} \uline{models} \uline{buil\over{d}{t} up}
\uline{with k-modules which are} \uuline{projective}, \uline{or at} any rate flat (namely torsi-%
on free, if $k=\underline{Z}$ ) .

\margin{14}
